\documentclass[10pt,a4paper]{amsart}
\usepackage[margin=19mm]{geometry}
\usepackage{amsmath,amssymb,mathtools}
\usepackage[scr=rsfs,cal=euler]{mathalfa}
\usepackage{xfrac}
\usepackage[unicode,hidelinks,bookmarks=false]{hyperref}
\newcommand{\ie}{i.e.\ }
\newcommand{\cf}{cf.\ }
\newcommand{\resp}{respectively\ }
\newcommand{\Subsection}{\subsection}
\providecommand{\nobreakdash}{\nobreak\hbox{-}}
\setlength{\emergencystretch}{2em}
\multlinegap0pt
\usepackage{amssymb}
\usepackage[shortlabels]{enumitem}
\usepackage{mathtools}
\usepackage{float}
\usepackage{multirow}
\usepackage{bbm}
\usepackage{xfrac}
\newtheorem{lemma}{Lemma}
\newcommand{\A}{\mathrm{A}}
\DeclareMathOperator{\Aut}{Aut}
\newcommand{\C}{\mathrm{C}}
\newcommand\FF{\mathbb{F}}
\DeclareMathOperator{\Miy}{Miy}
\newcommand\al{\alpha}
\newcommand\bt{\beta}
\DeclareMathOperator{\ch}{char}
\newcommand{\half}{{\tau_{\sfrac{1}{2}}}} % keep the extra brackets so we can use this as a sub/superscript
\newcommand\la{\langle}
\newcommand\lm{\lambda}
\renewcommand{\phi}{\varphi}
\newcommand\ra{\rangle}
\title{A guide to the $2$-generated axial algebras of Monster type}
\author{J.~M\textsuperscript{c}Inroy \and A.W.~Mir}
\date{Source: arXiv:2512.24987v1 (CC BY 4.0)}
\begin{document}
\maketitle

\noindent\emph{Source and licence.} A guide to the $2$-generated axial algebras of Monster type. Version arXiv:2512.24987v1 (CC BY 4.0), distributed by arXiv under the Creative Commons Attribution 4.0 International licence, http://creativecommons.org/licenses/by/4.0/, as recorded in the arXiv licence metadata for this exact version and shown on the abstract page https://arxiv.org/abs/2512.24987v1. Full licence notice: \texttt{licenses/arxiv-2512.24987-CC-BY-4.0.txt}.\par\smallskip

\noindent\textit{Editorial scope.} Counted unit: the article's lemma 4/7 (source0.tex lines 1830--1836). The article's complete original proof of it is delivered with it (source0.tex lines 1837--1867, one \texttt{proof} environment of 31 lines). No mathematics is written, repaired or re-derived: the statement and the proof are the article's own lines, in the article's own notation, and no retained line is substituted or re-flowed.  No further source-local context is needed: every cross-reference of every retained line resolves inside the two delivered spans.  Nothing was omitted from any retained span, so the package carries no omission record.  No retained line contains a citation, so the wrapper carries no bibliography and no bibliography entry.  No retained line contains a figure environment, an image-inclusion command or drawing code, so no image is delivered and the package declares no asset dependency.  Editorial formatting only, in addition: an AMS \texttt{amsart} preamble that restates the article's own declarations and macros used by the retained lines, namely the environments \texttt{lemma} on separate counters, together with the standard packages those lines need.  The retained environments print in this document as Lemma 1 in order of appearance, whereas the article prints the same items with the numbering of its own full document.  The complete native source of the article as distributed in the arXiv e-print for this exact version is bundled unchanged as \texttt{sources/191950/source0.tex}.
\begin{lemma}\label{6Aideal 4/7}
Suppose that $\ch \FF \neq 5$ and $\al = \frac{4}{7}$, then the radical of $A := 6\A(\frac{4}{7}, -\frac{4}{7})$ is
\[
R = \la a_0 - a_2, a_0 - a_{-2}, a_3 - a_{-1}, a_3 - a_1, z \ra
\]
and there are no non-trivial proper subideals of $R$.  The quotient $A/R$ is $3\C(\frac{4}{7})$.
\end{lemma}

\begin{proof}
Note that $\al = \frac{4}{7}$, the characteristic of the field is not $3$, or $7$.  In addition, the characteristic is not $11$ either as in this case $\bt = -\frac{4}{7} = 1$, a contradiction.  By computation, we find that the characteristic polynomial of the Frobenius form is $t^5(t - \frac{15}{7})(t^2 - \frac{34}{7}t + \frac{165}{49})$.  Since the characteristic is not $3$, $7$, or $11$ and is not $5$ by assumption, $15$ and $165 = 3\cdot 5 \cdot 11$ are not zero.  Hence the nullspace $R$ is at most $5$-dimensional in any characteristic other than $5$.  On the other hand, it is easy to see that $a_0 - a_2, a_0 - a_{-2}, a_3 - a_{-1}, a_3 - a_1, z \in R$ and hence the the radical is indeed $5$-dimensional.

Since $\Miy(A) \cong S_3$, $R$ decomposes as the direct sum of irreducible $S_3$-modules.  It is easy to see that it is the sum of a trivial module $\la z \ra \cong U_1$ and two $2$-dimensional irreducible modules $\langle a_0 - a_2, a_0 - a_{-2} \rangle, \langle a_3 - a_{-1}, a_3 - a_1 \rangle \cong U_2$.  Let $I \trianglelefteq A$ be a non-trivial proper subideal of $R$. The only possible $G$-mod structure for $I$ is $U_1$, $U_2$, $U_1 \oplus U_2$, or $U_2 \oplus U_2$.

First suppose that $z \in I$.  Then we claim that $I = R$.  Indeed,
\begin{equation}\label{6A 4/7 eqn z}
-\tfrac{7}{2} a_0 \cdot z = 2a_0 - a_{-2} - a_2 + z = (a_0 - a_2) + (a_0 -a_{-2}) + z
\end{equation}
and so, as $\ch \FF \neq 7$, $(a_0 - a_2) + (a_0 -a_{-2}) \in I$.  There exists $\phi \in \Aut(A)$, such that $\phi \colon a_i \mapsto a_{i+3}$ and fixes $z$; indeed this is $\phi = \half \tau_2$.  Now, by applying $\phi$ to both side above, we get $-\tfrac{7}{2} a_3 \cdot z = (a_3 -a_{-1}) + (a_3 - a_1) +z \in I$. Hence $I = R$.  Therefore, the trivial module $\la z \ra$ cannot be a constituent of any proper ideal $I$ and in particular, $I$ cannot be $1$-dimensional.

So now suppose that $I$ contains a $2$-dimensional module $\la u , v \ra$ isomorphic to $U_2$.  By considering the action of $\Miy(A) \cong S_3$ on $\langle a_0 - a_2, a_0 - a_{-2} \rangle, \langle a_3 - a_{-1}, a_3 - a_1 \rangle$, we see any isomorphism between these is a scalar multiple of the map $\phi$.  So, there exists $\lm, \mu \in \mathbb{F}$, with $\lm$ and $\mu$ not both zero, such that $u := \lambda(a_0 - a_2) + \mu(a_3 - a_{-1})$ and $v :=  \lambda(a_0 - a_{-2}) + \mu(a_3 - a_1)$.  Now
\begin{align*}
7a_0 u &= a_0 \big( \lambda(a_0 - a_2) + \mu(a_3 - a_{-1}) \big) \\
&=  7\lambda a_0 - \lambda \left(a_0 + a_2 + 5a_{-2} - 3z \right) + 2\mu \left(a_0 + a_3 - c\right) \\
&\phantom{{}={}} + 2\mu \left(a_0 + a_{-1} - a_1 - a_2 - a_{-2} + c + z\right) \\
&= (6\lm + 4 \mu)a_0 +2\mu (a_{-1} - a_1) - (\lm + 2\mu) a_2 \\ 
&\phantom{{}={}} - (5\lm + 2\mu) a_{-2} + (3\lm + 2\mu) z
\end{align*}
In particular, if $3\lambda + 2\mu \neq 0$, then $z \in I$, and so by the previous argument $I = R$.  Since $I$ is a proper ideal of $R$, we must have $3\lambda + 2\mu = 0$.
%\begin{align*}
%7a_1 u &= a_1 \left( \lambda(a_0 - a_2) + \mu(a_3 - a_{-1}) \right) \\
%&= 2\lambda\left(a_0 + a_3 - c\right) + 2\lm \left(a_2 + a_3 - a_{-2} - a_{-1}-a_1 -a_0 + c +z\right) \\
%&\phantom{{}={}} + 7\mu a_3   - \mu \left(a_3 - a_{-1}+ 5 a_1 -3 z\right) \\
%&=-\left(2\lm + 5\mu\right)a_1 + 2\lm (a_2 - a_{-2}) + \left(4\lm + 6\mu \right)a_3 \\
%&\phantom{{}={}} - \left(2\lm + \mu \right)a_{-1}  + \left(2\lm + 3\mu \right)z
%\end{align*}
Now, note that applying $\phi$ to the above maps $a_0$ to $a_3$ and switches the roles of $\lambda$ and $\mu$.  Hence we also require $3\mu + 2\lm = 0$.  We cannot simultaneously have $3\lm + 2\mu = 0$ and $3\mu + 2\lm = 0$ unless $\ch \FF = 5$ and $\lm = \mu$.  Therefore $R$ has no non-trivial proper subideals.

Note that the quotient $A/R$ is generated by two axes $\bar{a}_0 = \bar{a}_2 = \bar{a}_{-2}$ and $\bar{a}_1 = \bar{a}_3 = \bar{a}_{-1}$ and the $\bt$-eigenspace of $a_0$ is contained in $R$.  It is then easy to see that $A/R$ is $2$-generated of Jordan type $\frac{4}{7}$ and hence is isomorphic to $3\C(\frac{4}{7})$.
\end{proof}

\end{document}
